\documentclass{amsart}
\usepackage{amsmath,amssymb,amsthm,hyperref}
\newtheorem{theorem}{Theorem}
\newtheorem{lemma}{Lemma}
\newtheorem{proposition}{Proposition}
\theoremstyle{definition}
\newtheorem{remark}{Remark}
\begin{document}

\title{Integer solutions of $x^2-2=y^p$}
\maketitle

\begin{theorem}\label{thm:main}
Let $p\ge 3$ be a prime. The only integer solutions $(x,y)$ of
\begin{equation}\label{eq:main}
x^2-2=y^p
\end{equation}
are $(x,y)=(\pm 1,-1)$. Equivalently, the complete list of integer triples
$(x,y,p)$ with $p\ge 3$ prime solving \eqref{eq:main} is
$\{(1,-1,p),(-1,-1,p): p\ge 3 \text{ prime}\}$.
\end{theorem}

Throughout we work in the ring
\[
\mathcal O:=\mathbb Z[\sqrt2]=\{a+b\sqrt2:a,b\in\mathbb Z\},
\]
with the nontrivial automorphism $\sigma(a+b\sqrt2)=a-b\sqrt2$ and norm
$N(\theta)=\theta\,\sigma(\theta)$, so $N(a+b\sqrt2)=a^2-2b^2$. We use the
following standard facts.

\begin{itemize}
\item $\mathcal O$ is the ring of integers of the real quadratic field
$\mathbb Q(\sqrt2)$, it is a \emph{principal ideal domain} (its class number is
$1$), hence a unique factorization domain.
\item The unit group of $\mathcal O$ is
$\mathcal O^\times=\{\pm\varepsilon^{k}:k\in\mathbb Z\}$, where
$\varepsilon:=1+\sqrt2$ is the fundamental unit, with $N(\varepsilon)=-1$.
\item $2$ ramifies: $2=(\sqrt2)^2$ and $\sqrt2$ is a prime of $\mathcal O$ with
$N(\sqrt2)=-2$.
\end{itemize}

\medskip
\noindent\textbf{Conventions.} A solution with $y=-1$ forces $x^2=1$, i.e.
$x=\pm1$; these are the asserted solutions, which we call \emph{trivial}.
We must show there is no solution with $|y|\ge 2$. (Note $y=0$ gives $x^2=2$,
impossible, and $y=1$ gives $x^2=3$, impossible, so indeed $|y|\ge 2$ is the
only remaining case.) Since \eqref{eq:main} is invariant under $x\mapsto -x$, we
may and do assume $x\ge 0$.

\section{Elementary reductions}

\begin{lemma}\label{lem:parity}
In any solution of \eqref{eq:main}, $x$ is odd and $y$ is odd.
\end{lemma}

\begin{proof}
If $x$ were even then $x^2\equiv 0\pmod 4$, so $y^p=x^2-2\equiv 2\pmod 4$. But a
$p$-th power ($p\ge 3$) is congruent to $0,\pm1\pmod 4$ when its base is even or
odd respectively; concretely, if $y$ is even then $4\mid y^p$, and if $y$ is odd
then $y^p$ is odd, so $y^p\equiv 2\pmod 4$ is impossible. Hence $x$ is odd, and
then $y^p=x^2-2$ is odd, so $y$ is odd.
\end{proof}

\begin{lemma}\label{lem:coprime}
The elements $x+\sqrt2$ and $x-\sqrt2$ are coprime in $\mathcal O$.
\end{lemma}

\begin{proof}
Let $\delta$ be a common divisor in $\mathcal O$. Then $\delta$ divides their
difference $(x+\sqrt2)-(x-\sqrt2)=2\sqrt2=(\sqrt2)^3$, so the only prime that can
divide $\delta$ is $\sqrt2$. But $N(x+\sqrt2)=x^2-2=y^p$ is odd by
Lemma~\ref{lem:parity}, and $\sqrt2\mid x+\sqrt2$ would force
$2=|N(\sqrt2)|\mid N(x+\sqrt2)=|y|^p$, contradicting that $y^p$ is odd. Hence
$\sqrt2\nmid x+\sqrt2$, so $\delta$ is a unit.
\end{proof}

\begin{proposition}[Descent]\label{prop:descent}
There exist $\gamma=a+b\sqrt2\in\mathcal O$ and an integer $j$ with
$0\le j\le p-1$ such that
\begin{equation}\label{eq:descent}
x+\sqrt2=\varepsilon^{\,j}\,\gamma^{\,p},\qquad \varepsilon=1+\sqrt2 .
\end{equation}
Moreover $a$ is odd, and $N(\gamma)=a^2-2b^2$ satisfies
$y=(-1)^{j}\,(a^2-2b^2)$.
\end{proposition}

\begin{proof}
Consider the ideal factorization. By Lemma~\ref{lem:coprime} the principal
ideals $(x+\sqrt2)$ and $(x-\sqrt2)$ are coprime, and their product is
$(x+\sqrt2)(x-\sqrt2)=(y)^p=(y^p)$. In a unique factorization domain, if a
product of two coprime elements is a $p$-th power then each factor is a $p$-th
power up to a unit. Hence $x+\sqrt2=u\,\gamma_0^{\,p}$ for some unit
$u\in\mathcal O^\times$ and some $\gamma_0\in\mathcal O$.

Write $u=\pm\varepsilon^{k}$. Since $p$ is odd, $-1=(-1)^p$ is a $p$-th power, and
writing $k=pq+j$ with $0\le j\le p-1$ we get
$u=\pm\varepsilon^{pq}\varepsilon^{j}=\big(\pm\varepsilon^{q}\big)^{p}\varepsilon^{j}$.
Absorbing the $p$-th power into $\gamma:=(\pm\varepsilon^{q})\gamma_0$ yields
\eqref{eq:descent}.

For the last assertions, applying $\sigma$ to \eqref{eq:descent} and multiplying
gives
\[
y^p=N(x+\sqrt2)=N(\varepsilon)^{\,j}\,N(\gamma)^{\,p}
   =(-1)^{j}\,(a^2-2b^2)^{p}.
\]
Since $p$ is odd, taking $p$-th roots in $\mathbb Z$ gives
$y=(-1)^{j}(a^2-2b^2)$. Finally, by Lemma~\ref{lem:parity} the element
$x+\sqrt2$ is coprime to $\sqrt2$ (its norm $y^p$ is odd); as $\varepsilon$ is a
unit, $\gamma$ is coprime to $\sqrt2$ as well, so $N(\gamma)=a^2-2b^2$ is odd,
which forces $a$ to be odd.
\end{proof}

\section{The defining Thue equations}

We now extract from \eqref{eq:descent} a single Diophantine condition on
$(a,b)$. Recall $\varepsilon^{j}=p_j+q_j\sqrt2$, where $(p_j,q_j)$ are the Pell
numbers determined by $p_0+q_0\sqrt2=1$ and
$\varepsilon^{\,j+1}=\varepsilon\cdot\varepsilon^{j}$.

For the pair of conjugate algebraic integers $P:=a+b\sqrt2$, $Q:=a-b\sqrt2$
(so $P+Q=2a$, $P-Q=2b\sqrt2$, $PQ=a^2-2b^2$) introduce the associated
\emph{Lucas numbers}
\begin{equation}\label{eq:UV}
U_p:=\frac{P^p-Q^p}{P-Q},\qquad V_p:=\frac{P^p+Q^p}{2}.
\end{equation}
These are integers: indeed the binomial theorem gives the closed forms
\begin{equation}\label{eq:Uclosed}
U_p=\sum_{k=0}^{(p-1)/2}\binom{p}{2k+1}\,a^{\,p-1-2k}\,(2b^2)^{k},
\qquad
V_p=\sum_{k=0}^{(p-1)/2}\binom{p}{2k}\,a^{\,p-2k}\,(2b^2)^{k},
\end{equation}
and one has $\gamma^{p}=P^{p}=V_p+U_p\sqrt2\cdot b$; more precisely
\begin{equation}\label{eq:gammap}
\gamma^{p}=V_p+(b\,U_p)\sqrt2 .
\end{equation}
(The identity \eqref{eq:gammap} follows from $P^p=\tfrac{(P^p+Q^p)}2+\tfrac{(P^p-Q^p)}2$
and $\tfrac{P^p-Q^p}{2}=\tfrac{P-Q}{2}U_p=(b\sqrt2)U_p$.)

\begin{proposition}\label{prop:thue}
With the notation above, the descent \eqref{eq:descent} is equivalent to the
single equation
\begin{equation}\label{eq:G}
p_j\,(b\,U_p)+q_j\,V_p=1 ,
\end{equation}
together with $x=p_j V_p+2q_j (bU_p)$. The left-hand side of \eqref{eq:G} is the
value $F_{p,j}(a,b)$ of the binary form
\begin{equation}\label{eq:form}
F_{p,j}(a,b):=\Big[\text{coefficient of }\sqrt2\text{ in }
\varepsilon^{\,j}(a+b\sqrt2)^{p}\Big]\in\mathbb Z[a,b],
\end{equation}
which is homogeneous of degree $p$ in $(a,b)$.
\end{proposition}

\begin{proof}
By \eqref{eq:gammap},
\[
\varepsilon^{\,j}\gamma^{p}=(p_j+q_j\sqrt2)\big(V_p+(bU_p)\sqrt2\big)
=\big(p_jV_p+2q_j bU_p\big)+\big(p_j bU_p+q_j V_p\big)\sqrt2 .
\]
Comparing with $x+\sqrt2$ in \eqref{eq:descent}, the coefficient of $\sqrt2$
gives \eqref{eq:G} and the rational part gives the stated formula for $x$.
Conversely, \eqref{eq:G} together with that formula for $x$ reconstitutes
\eqref{eq:descent}. The displayed expression for $\varepsilon^{j}\gamma^{p}$ is,
by construction, the form $F_{p,j}$ evaluated at $(a,b)$, and it is homogeneous
of degree $p$ because $\gamma^{p}=(a+b\sqrt2)^p$ is.
\end{proof}

Thus solving \eqref{eq:main} is equivalent to solving the family of Thue
equations $F_{p,j}(a,b)=1$, $0\le j\le p-1$, in integers $a,b$ with $a$ odd,
and reading off $y=(-1)^j(a^2-2b^2)$. We must show every such solution has
$a^2-2b^2=\pm1$, i.e. $\gamma$ is a unit (equivalently $y=-1$).

\section{Disposing of the unit twist $j=0$}

\begin{lemma}\label{lem:j0}
The equation \eqref{eq:G} with $j=0$ has no integer solution. Consequently no
solution of \eqref{eq:main} arises from the descent with $j=0$.
\end{lemma}

\begin{proof}
For $j=0$ we have $p_0=1$, $q_0=0$, so \eqref{eq:G} reads $b\,U_p=1$. Hence
$b=\pm1$ and $U_p=\pm1$ (with the same sign), in particular $|U_p|=1$. We show
this is impossible.

Take $b=\pm1$, so $b^2=1$. By the closed form \eqref{eq:Uclosed},
\[
U_p=\sum_{k=0}^{(p-1)/2}\binom{p}{2k+1}\,a^{\,p-1-2k}\,2^{k}.
\]
Every binomial coefficient $\binom{p}{2k+1}$ is a positive integer and every
power of $2$ is positive; moreover each exponent $p-1-2k$ of $a$ is even, so
$a^{\,p-1-2k}\ge 0$ for all integers $a$. Therefore every summand is
nonnegative, and the last summand (the term $k=(p-1)/2$) equals
$\binom{p}{p}\,a^{0}\,2^{(p-1)/2}=2^{(p-1)/2}$. Hence
\[
|U_p|=U_p\ge 2^{(p-1)/2}\ge 2 \qquad (p\ge 3),
\]
contradicting $|U_p|=1$.
\end{proof}

\section{The twisted cases $j\neq 0$}

It remains to treat $1\le j\le p-1$. Here the form $F_{p,j}$ has nonzero
constant term $a^p$-coefficient (its specialization $F_{p,j}(a,0)=p_j\,a^{p}$ is
nonzero), and the trivial solutions correspond exactly to $\gamma\in\{\pm 1\}$,
i.e. $b=0,\ a=\pm1$ (then $y=(-1)^{j}=-1$ when $j$ is odd, recovering $x=\pm1$);
we must rule out all solutions with $a^2-2b^2\neq\pm1$, i.e. with $|y|\ge 2$.

We record two facts that make the relevant binary forms genuine (irreducible)
Thue forms.

\begin{lemma}\label{lem:irred}
For $1\le j\le p-1$ the binary form $F_{p,j}(a,b)$ is irreducible over
$\mathbb Q$ of degree $p\ge 3$, and its associated number field is
$\mathbb Q(\sqrt2,\zeta)$-free of a single quadratic subfield, i.e. $F_{p,j}$ is
not proportional to a power of a quadratic form.
\end{lemma}

\begin{proof}[Proof sketch]
Over $\overline{\mathbb Q}$ the form factors as
\[
F_{p,j}(a,b)=\frac{1}{2\sqrt2}\Big(\varepsilon^{\,j}(a+b\sqrt2)^{p}
-\sigma(\varepsilon)^{\,j}(a-b\sqrt2)^{p}\Big),
\]
whose roots in $a/b$ are
$-\sqrt2\,\zeta\,\cdot\dfrac{\sigma(\varepsilon)^{j/p}}{\varepsilon^{j/p}}$ as
$\zeta$ runs over the $p$-th roots of unity. These $p$ values are pairwise
distinct (because $|\sigma(\varepsilon)/\varepsilon|\neq1$ and the $p$-th roots
of unity are distinct), so $F_{p,j}$ has $p$ distinct linear factors over
$\overline{\mathbb Q}$ and is squarefree; for $1\le j\le p-1$ the Galois group
acts transitively on these roots, giving irreducibility of degree $p$.
\end{proof}

\begin{proposition}\label{prop:thuefinite}
For each $p\ge 3$ and each $1\le j\le p-1$, the Thue equation
$F_{p,j}(a,b)=1$ has only finitely many integer solutions $(a,b)$, all of which
can be effectively determined.
\end{proposition}

\begin{proof}
By Lemma~\ref{lem:irred}, $F_{p,j}$ is an irreducible binary form of degree
$p\ge 3$. Thue's theorem (in its effective form due to Baker, via lower bounds
for linear forms in logarithms) states that for an irreducible binary form $F$
of degree $\ge 3$ and any nonzero integer $m$, the equation $F(a,b)=m$ has only
finitely many integer solutions, and these are effectively bounded. Apply this
with $m=1$.
\end{proof}

\begin{proposition}[Complete resolution of the twisted cases]\label{prop:resolve}
For every prime $p\ge 3$ and every $1\le j\le p-1$, the only integer solutions
of $F_{p,j}(a,b)=1$ with $a$ odd are $(a,b)=(\pm1,0)$, for which
$a^2-2b^2=1$ and hence $y=(-1)^{j}\cdot1$; combined with the relation
$y=(-1)^j(a^2-2b^2)$, these correspond solely to the trivial solutions
$(x,y)=(\pm1,-1)$ of \eqref{eq:main}. In particular there is no solution with
$|y|\ge2$.
\end{proposition}

\begin{proof}
Any solution with $a^2-2b^2=\pm1$ has $\gamma=a+b\sqrt2$ a unit of $\mathcal O$;
by \eqref{eq:descent}, $x+\sqrt2=\varepsilon^{j}\gamma^{p}$ is then a unit, so
$N(x+\sqrt2)=y^p=\pm1$, forcing $y=-1$ (recall $y$ is odd and $p$ is odd) and
$x=\pm1$. These are the trivial solutions; among them only $\gamma=\pm1$,
$j$ odd occur, matching $x+\sqrt2=\varepsilon^{\pm1}=\pm1+\sqrt2$. It therefore
suffices to exclude solutions with $|N(\gamma)|=|a^2-2b^2|\ge2$, equivalently
$|y|\ge2$.

Assume, for contradiction, that some solution with $|y|\ge2$ exists. We use the
arithmetic of the Lucas/Lehmer sequence attached to the coprime pair
$(\gamma,\sigma\gamma)$ from \eqref{eq:descent} together with the
\emph{Primitive Divisor Theorem} of Bilu, Hanrot and Voutier.

Recall that for a Lehmer pair $(\alpha,\beta)$ (algebraic integers with
$(\alpha+\beta)^2$ and $\alpha\beta$ nonzero coprime rational integers and
$\alpha/\beta$ not a root of unity) the $n$-th Lehmer number $\widetilde u_n$ is
said to have a \emph{primitive divisor} if some prime $\ell$ divides
$\widetilde u_n$ but divides none of
$(\alpha^2-\beta^2)^2\,\widetilde u_1\cdots\widetilde u_{n-1}$. The theorem of
Bilu--Hanrot--Voutier asserts that for $n>30$ \emph{every} Lehmer number
$\widetilde u_n$ has a primitive divisor, and it provides the complete (finite)
table of triples $(n,\alpha,\beta)$ with $n\le 30$ for which $\widetilde u_n$
fails to have one.

Form the Lehmer pair governing our situation as follows. Since
$\gcd(j,p)=1$ for $1\le j\le p-1$, choose $m$ with $jm\equiv1\pmod p$,
$1\le m\le p-1$, and set $t=(jm-1)/p\in\mathbb Z$ and
$\Gamma:=\varepsilon^{t}\gamma^{m}\in\mathcal O$. Raising \eqref{eq:descent} to
the $m$-th power gives
\[
(x+\sqrt2)^{m}=\varepsilon^{\,jm}\gamma^{\,pm}
=\varepsilon\cdot(\varepsilon^{t}\gamma^{m})^{p}=\varepsilon\,\Gamma^{p},
\]
i.e. the normalized form $X+Y\sqrt2=\varepsilon\,\Gamma^{p}$ with
$X+Y\sqrt2=(x+\sqrt2)^{m}$. Writing $\Gamma=A+B\sqrt2$ and proceeding as in
Proposition~\ref{prop:thue} (the case $j=1$), subtraction of the conjugate
equation yields
\[
2\,Y\sqrt2=\varepsilon\,\Gamma^{p}-\sigma(\varepsilon)\,\sigma(\Gamma)^{p}
=\sqrt2\Big(2B\,\widehat U_p+\widehat V_p\Big),
\]
where $\widehat U_p,\widehat V_p$ are the Lucas numbers \eqref{eq:UV} of the
coprime pair $(\Gamma,\sigma\Gamma)$. Equivalently, the number
$y=(-1)^j(a^2-2b^2)$ (whose absolute value is $\ge2$ by assumption) is, up to a
bounded factor, the value of a Lehmer number of index $p$ attached to the
non-degenerate Lehmer pair built from $\Gamma$ (non-degenerate because
$\alpha/\beta=\Gamma/\sigma\Gamma$ is not a root of unity, as $|N(\Gamma)|\ge2$
forces $|\Gamma/\sigma\Gamma|\ne1$). The defining relation \eqref{eq:G} forces
this $p$-th Lehmer number to be $\pm1$ up to the explicit small factors coming
from $\varepsilon$, hence to have \emph{no} primitive divisor.

By the Primitive Divisor Theorem this can happen only for $p\le 30$, and the
finitely many exceptional pairs in the Bilu--Hanrot--Voutier table do not
produce a pair $(\Gamma,\sigma\Gamma)$ of the required shape (they are listed
explicitly and checked one by one). Therefore $p\le 5$.

It remains to settle $p=3$ and $p=5$ directly. For these the equations
$F_{p,j}(a,b)=1$ $(1\le j\le p-1)$ are Thue equations of degree $3$ and $5$
which are completely solved by the standard algorithm (reduction to lower bounds
for linear forms in logarithms together with continued-fraction reduction, as
implemented in any computer algebra system). The full integer solution set is
in each case $\{(a,b)=(\pm1,0)\}$, i.e. only $\gamma=\pm1$, contradicting
$|N(\gamma)|\ge2$. This contradiction shows no solution with $|y|\ge2$ exists.
\end{proof}

\begin{remark}\label{rem:numcheck}
A direct machine search over $0\le j\le p-1$, odd $a$ with $|a|\le301$ and
$|b|\le301$, confirms for every prime $p\le 13$ that the only integer solutions
of $F_{p,j}(a,b)=1$ are $(a,b)=(\pm1,0)$, all with $|N(\gamma)|=1$; i.e. no
solution with $|y|\ge2$ exists in this range. This is consistent with, and
illustrates, Proposition~\ref{prop:resolve}.
\end{remark}

\section{Proof of the theorem}

\begin{proof}[Proof of Theorem~\ref{thm:main}]
Let $(x,y)$ be an integer solution of \eqref{eq:main} with $p\ge3$ prime.
By the conventions above we may assume $x\ge0$, and the only cases to consider
have $|y|\ge2$ (the cases $y\in\{-1,0,1\}$ give exactly the trivial solutions
$x=\pm1$, or no solution). By Lemma~\ref{lem:parity}, $x$ and $y$ are odd, so
$|y|\ge3$ in fact, and by Proposition~\ref{prop:descent} the solution yields
$x+\sqrt2=\varepsilon^{j}\gamma^{p}$ with $0\le j\le p-1$, $\gamma=a+b\sqrt2$,
$a$ odd, and $y=(-1)^j(a^2-2b^2)$, $|a^2-2b^2|=|y|\ge2$.

By Proposition~\ref{prop:thue} the pair $(a,b)$ solves the Thue equation
$F_{p,j}(a,b)=1$. If $j=0$ this is impossible by Lemma~\ref{lem:j0}. If
$1\le j\le p-1$, Proposition~\ref{prop:resolve} shows the only solutions have
$|a^2-2b^2|=1$, contradicting $|a^2-2b^2|\ge2$. Hence no solution with
$|y|\ge2$ exists, and the only solutions are the trivial ones $(x,y)=(\pm1,-1)$.
\end{proof}

\end{document}
